\chapter{Bridges from vectors to neural population geometry}\label{ch:bridges-linear}
\begin{readergoal}Derive basic linear-algebra tools from definitions and see exactly how they lead to PCA, projections, and regression.\end{readergoal}

\bridgeentry{span-basis}{Linear independence makes coordinate representations unique}
\begin{proposition}If $v_1,\ldots,v_k$ are linearly independent vectors, every vector in their span has exactly one representation $x=\sum_{j=1}^k a_jv_j$.\end{proposition}
\begin{proof}Existence follows from the definition of span. If $x=\sum a_jv_j=\sum b_jv_j$, subtracting gives $\sum(a_j-b_j)v_j=0$. Linear independence forces $a_j-b_j=0$ for each $j$, so coefficients are unique.\end{proof}
\paragraph{Example.} A latent subspace can have unique coordinates after fixing a basis, but another rotated basis describes the same subspace. ``Component 1'' need not be a unique biological source.

\bridgeentry{rank-nullity}{Dimensions split into visible and invisible directions}
\begin{theorem}[Rank--nullity]For a linear map $A:V\to W$ with $\dim V=n<\infty$, $\dim\ker A+\dim\operatorname{im}A=n$.\end{theorem}
\begin{proof}Choose a basis $u_1,\ldots,u_r$ of $\ker A$ and extend it to a basis of $V$ with $v_1,\ldots,v_{n-r}$. Their images $Av_1,\ldots,Av_{n-r}$ span $\operatorname{im}A$. If a linear combination of these images vanishes, the corresponding combination of $v_j$ lies in the kernel and must have zero coefficients by independence of the full basis. Thus the images form a basis of dimension $n-r$.\end{proof}
\paragraph{Example.} A lead-field map with nonzero nullspace cannot uniquely recover all neural current distributions from extracellular electrodes.

\bridgeentry{inner-cs}{Cauchy--Schwarz bounds the inner product}
\begin{theorem}[Cauchy--Schwarz]For a real inner-product space, $|\langle x,y\rangle|\leq\|x\|\|y\|$.\end{theorem}
\begin{proof}For $y\neq0$, nonnegativity of $\|x-\alpha y\|^2$ with $\alpha=\langle x,y\rangle/\|y\|^2$ gives $\|x\|^2-\langle x,y\rangle^2/\|y\|^2\geq0$. Rearranging and taking square roots proves the claim. If $y=0$ both sides vanish.\end{proof}
\paragraph{Example.} Applying this inequality to centered random variables implies $|\rho|\leq1$ when their variances are positive.

\bridgeentry{norm-triangle}{The triangle inequality follows from Cauchy--Schwarz}
\begin{proposition}For the norm $\|x\|=\sqrt{\langle x,x\rangle}$, $\|x+y\|\leq\|x\|+\|y\|$.\end{proposition}
\begin{proof}Expand $\|x+y\|^2=\|x\|^2+2\langle x,y\rangle+\|y\|^2$. The previous bridge bounds the middle term by $2\|x\|\|y\|$, making the sum no larger than $(\|x\|+\|y\|)^2$. Nonnegative square roots preserve the inequality.\end{proof}
\paragraph{Example.} Reconstruction error from two perturbations is at most the sum of their norm bounds; that inequality says nothing about whether the model is physiologically meaningful.

\bridgeentry{pythagoras}{Orthogonal components add energy, not amplitude}
\begin{proposition}If $\langle x,y\rangle=0$, then $\|x+y\|^2=\|x\|^2+\|y\|^2$.\end{proposition}
\begin{proof}The general squared-norm expansion has the cross term $2\langle x,y\rangle$. Orthogonality makes this term zero.\end{proof}
\paragraph{Example.} Orthogonal Fourier components add squared energy over a complete ideal window. Raw amplitudes are not generally additive as magnitudes.

\bridgeentry{gram-schmidt}{Independent vectors can be orthogonalized}
\begin{theorem}[Gram--Schmidt procedure]Every finite linearly independent family $v_1,\ldots,v_k$ in an inner-product space spans an orthonormal family $q_1,\ldots,q_k$ with matching partial spans.\end{theorem}
\begin{proof}Set $u_1=v_1$ and $q_1=u_1/\|u_1\|$. Recursively let $u_j=v_j-\sum_{i<j}\langle v_j,q_i\rangle q_i$. This vector is perpendicular to preceding $q_i$ by cancellation of inner products, and it cannot be zero without making $v_j$ a combination of preceding vectors. Normalize $q_j=u_j/\|u_j\|$ and proceed inductively.\end{proof}
\paragraph{Example.} Orthonormal bases permit stable projections, but finite-precision implementations often prefer Householder QR over classical Gram--Schmidt.

\bridgeentry{orthogonal-projection}{Projection solves constrained least squares}
\begin{proposition}If columns of $Q\in\mathbb R^{n\times k}$ are orthonormal, the closest vector to $y$ in $\operatorname{im}Q$ is $QQ^\top y$.\end{proposition}
\begin{proof}Write $p=QQ^\top y$ and $r=y-p$. Then $Q^\top r=Q^\top y-Q^\top QQ^\top y=0$. For any $z\in\operatorname{im}Q$, $y-z=r+(p-z)$ is a sum of orthogonal vectors, so $\|y-z\|^2=\|r\|^2+\|p-z\|^2\geq\|r\|^2$. The minimum occurs at $z=p$.\end{proof}
\paragraph{Example.} Linear neural decoding fits a projection onto a predictor subspace. Validation must still separate training and held-out trials.

\bridgeentry{covariance-psd}{Covariance matrices are positive semidefinite}
\begin{proposition}For a random vector $X\in\mathbb R^n$ with finite second moments, its covariance matrix $C$ is symmetric positive semidefinite: $v^\top Cv\geq0$ for all $v$.\end{proposition}
\begin{proof}Symmetry follows from $\Cov(X_i,X_j)=\Cov(X_j,X_i)$. Expanding the variance of the scalar $v^\top X$ yields $v^\top Cv=\Var(v^\top X)\geq0$.\end{proof}
\paragraph{Example.} PCA uses the eigendecomposition of $C$ and nonnegative eigenvalues. A numerical covariance estimate with materially negative eigenvalues calls for an implementation or modeling audit.

\begin{exercises}
\item Find a basis and nullspace for $A=\begin{bmatrix}1&1\\2&2\end{bmatrix}$.
\item Derive the ordinary least-squares normal equations by orthogonality.
\item Explain why PCA axes can change substantially when leading eigenvalues are nearly equal.
\end{exercises}